\documentclass[11pt]{article}
\usepackage{amsmath,amssymb,amsthm,geometry}
\geometry{margin=.9in}\setlength{\emergencystretch}{1.5em}
\newtheorem{theorem}{Theorem}\newtheorem{lemma}{Lemma}
\newcommand{\dd}{\,d}\newcommand{\E}{\mathbb E}
\title{Cancellation of the transverse source and the critical weak clock}
\author{E5 closure attempt, step 7}\date{Version 1}
\begin{document}\maketitle

\paragraph{Question.}
Use the actual isotropic aligned balanced critical family
$S_0=cq^{1/\omega}$ with fixed $c>0$, $\lambda=\rho^2$,
$\omega=(1-\lambda)/2$, $L=\lambda/2$ and WIT's witness rectangle.
Let $s=t-t_h$ be time after first right half mass. From a fixed COD
witness $t_0=t_h+s_0$, set $u=r_{t_0}(\xi)/|r_{t_0}(\xi)|$,
$v=(-u_2,u_1)$, $r_t(z)=z-f_t(z)$,
$\xi=(\sin(q\zeta),-\cos(q\zeta))$. Write
$A=v^T(f_t-f_{t_0})(\xi)/q$ and $B=u^T(f_t-f_{t_0})(\xi)/q^2$;
$e_q=(|r_t(\xi)|^2-1)/(2q^2)$, $e_*=\min_J\mathcal E_{\rm coh}$.
The target is a trained time in the measured half-weak-response
window satisfying DRR's two strong/weak learning and loss tests and
$e_q<e_*$. We prove statements about the incoming trajectory,
then test their continuation to this endpoint.

\section{Net transverse growth and its signed sources}
The BRW extracted law at each finite $s$ has logistic mass
$M=(1+e^{-s})^{-1}$, zero weak mass, and fixed normalized rank
weights $4\cos^2\alpha\,d\alpha/(2\pi)$ on the right half.
Q4 is the bottom atom, Q1 is co-ordered in its own field.
Use $P=\Phi(\rho x)$, $k=K_s=\int K(\rho x)\dd\mu$,
$Y=\int y\dd\mu$, and $S_x=\int F(\rho x,\rho\tilde x)\dd\mu$.
Here $K(z)=\phi(z)+z\Phi(z)$, and $F(t,r)=\E[(Z+r)_+
{\bf1}_{Z>-t}]$. Put $H=\int(\zeta-x)_+\dd\mu$ and
$\mathcal A=\{x<\zeta\}$. Constants may depend on the fixed family
and witness but not on the extracted law.

\begin{lemma}[Location of the bulk]\label{an04tbwv1:bulk}
There are constants $b,d,C>0$ such that for all sufficiently large
$s$, except for normalized upper mass at most
$F_s=C\log(1+s)/(1+s)$,
\[
 -\sqrt{s-\tfrac{2}{3L}\log(1+s)+d}\le x(s)
       \le-\sqrt{s-b\log(1+s)} .
\]
The lower bound holds for all ranks. Consequently
\begin{align}
 H(s)&=\sqrt s+\zeta+O(\log(1+s)/\sqrt{s}),\nonumber\\
 \int [P+\phi(\rho x)]\dd\mu&=O(\log(1+s)/(1+s)).
 \label{an04tbwv1:exposure}
\end{align}
\end{lemma}
\begin{proof}
The first face is SPM's Gaussian lower face for the bottom path,
extended to all ranks by order. Its guards are precisely the BRW
logistic entrance: $2x_a'\ge-Ce^{Ls}(1+s)^{-2/3}P_{x_a}$ at
$x_a<0$, from the true donor and core bounds. For the stated
radicand $h$ with $d$ large, entry is strict at a fixed sufficiently
large $s_1$, $h'\ge1/2$, and $\Phi(-\rho\sqrt h)\le
\phi(\rho\sqrt h)/(\rho\sqrt h)$ gives $2x_a'>-h'/\sqrt h$ on
the sole lower face.

For the upper face fix the terminal $s$ and a rank $f\ge F_s$
(Q4 may use $f=1/2$). The BRW rank ceiling on $r\in[s-2,s]$ is
$C(1+f^{-1/3})e^{Lr-\lambda f r/16}$ for $x_++y_+$.
Choose the constant in $F_s$ large enough that this is
$o(D(s))$ uniformly, $D(s)=e^{Ls}/(1+s)$. For example require
$\lambda C/32>3$. At each $r$ the protected donors above
$d_1e^{Lr}$ have first moment at least $c_1D(r)$ with actual
mass; both conclusions follow from BRW's two-face moving barrier
for its logistic law. Those donors are above this moving receiver.
Thus $S_x\ge\rho P c_1D(r)$ even when $x<0$. The input equation
yields, whenever $x\ge-z_s$, $z_s=\sqrt{s-b\log(1+s)}$,
\[
 2x'\le e^{-(s-2)}z_s+\rho\phi(0)-c_2 D(s)P
\]
for large $s$. On $x\ge0$ the right side is at most $-c_3D(s)$.
Entry height is $o(D(s))$, so by $s-1$ the input is nonpositive
and cannot recross zero. On $-z_s\le x\le0$,
$P\ge c e^{-Ls}(1+s)^{Lb-1/2}$ by the lower Gaussian tail.
Take $Lb>2$. The negative speed dominates $\sqrt s$ here, so the
second unit interval suffices to reach $-z_s$. Both crossing faces
$0,-z_s$ forbid upward return during the corresponding intervals;
a rank already below a face requires no crossing time.
The positive radicand and all these bounds are only used for large $s$.
With $M=1-O(e^{-s})$, bounding the exposure above by the all-rank
lower face and below by the contribution of mass $M(1-F_s)$ proves
the first formula. On this bulk both Gaussian functions in the
second formula are exponentially small up to a fixed polynomial,
and on the remaining mass both are bounded.
\end{proof}

\begin{lemma}[Cost of overshooting the cross regression]\label{an04tbwv1:sign}
Uniformly for large $s$, $k'(s)>0$ and every active label satisfies
$y\le k(s)/\rho+C$.
\end{lemma}
\begin{proof}
First compute the centered overlap. For any $t,r$,
$0\le F(t,r)-K(r)\Phi(t)\le\phi(t)$: the lower sign follows from
covariance of increasing functions of the same normal. For the
upper bound subtract $(Z+r)_+$ from $Z$; the difference is
nondecreasing too, and $\operatorname{Cov}(Z,{\bf1}_{Z>-t})=\phi(t)$.
Write $S_x=kP+C_x$ with $0\le C_x\le M\phi(\rho x)$.

COD's active output ceiling and its positive-rank mass cutoff give
$\int_{\mathcal A}y_+\dd\mu=o(D(s))$: multiply
$O(e^{Ls}(1+s)^{-2/3})$ by $F_s$, noting the core output is negative
late. Let $p=\Phi(\rho\zeta)$, $w=\max(P^2,p^2)$. Order in the
single population gives $\int yw\dd\mu\ge Y\int w\dd\mu/M\ge p^2Y$
as $Y>0$. Replacing $w$ by $P^2$ costs at most
$p^2\int_{\mathcal A}y_+\dd\mu$. COD gives $Y\ge c_4 D(s)$.
Hence $\int yP^2\dd\mu\ge p^2Y/2$ late. In
$k'=(1-M)k+\rho\int P x'\dd\mu$, insert BRW field. By the
centered overlap inequality the doubled integral is bounded below by
\[
 2\int P x'\dd\mu\ge
 -(1-M)\int x_+\dd\mu-\rho k\int P^2\dd\mu+\lambda\int yP^2\dd\mu.
\]
BRW and COD bound $k$ and $\int x_+$ by
$C e^{Ls}(1+s)^{-2/3}$, while \eqref{an04tbwv1:exposure} gives
$\int P^2=O(F_s)$. The two costs are $o(Y)$.
Differentiation is justified by finite-interval first-moment domination
and $|K'|\le1$. This proves eventual increase.

Start at a fixed late time with $h=k/\rho+C$ strictly above all
$y$ for $x\le\zeta$, by COD's active-ceiling proof, including at
the gate (or directly its strict vertical-face test there).
Choose $C$ larger if necessary. At $x=\zeta,\ y\ge h$,
$2x'\ge-(1-M)\zeta+\rho(1-M)\phi(\rho\zeta)+\lambda p C>0$.
At $y=h,\ x\le\zeta$, $2y'\le\rho K(\rho\zeta)-Y<0<2h'$.
Both inequalities apply at the corner. These two faces exclude
entry into the forbidden open quadrant, proving the bound.
\end{proof}

\paragraph{Signed balance.}
Write $\ell_0=\zeta-H(s_0)$. In units of scaled transverse motion
the limiting positive cross source and the two signs of the mixed
weak term have densities on $\mathcal A$ (per actual radial mass)
\[
 c_x=\tfrac{\lambda}{2}(k/\rho-\ell_0)_+P,\qquad
 n_x=\tfrac{\lambda}{2}(y-\ell_0)_+P,\qquad
 p_x=\tfrac{\lambda}{2}(\ell_0-y)_+P .
\]
The signed transverse terms at late times are $+c_x,-n_x,+p_x$.
Since $k/\rho\ge c_3D(s)$ whereas $|\ell_0|$ is bounded,
the positive-part truncation of the cross coefficient is immaterial
late. Lemma~\ref{an04tbwv1:sign} gives $(n_x-c_x)_+\le C P$
on the active set. The actual exposure derivative, almost everywhere,
satisfies
\begin{align}
 H'={}&\int_{\mathcal A}(c_x-n_x+p_x)\dd\mu+R,\nonumber\\
 2R={}&(1-M)(H+\zeta\mu(\mathcal A))
              +\rho\int_{\mathcal A}(C_x-\phi(\rho x))\dd\mu .
 \label{an04tbwv1:balance}
\end{align}
Indeed differentiate the positive part along each path with the
common radial reaction, then insert the exact $x'$ and $S_x$.
This retains the logistic and Gaussian density terms.
By \eqref{an04tbwv1:exposure}, $\int_{s_1}^S |R|\le C(1+\log^2(1+S))$.
The same integrated bound holds for $\int_{\mathcal A}(n_x-c_x)_+\dd\mu$.
Consequently
\begin{equation}\label{an04tbwv1:matched}
 \int_{s_0}^S\int_{\mathcal A}\big[(c_x-n_x)_++p_x\big]\dd\mu\,ds
       =\sqrt S+O(1+\log^2(1+S)).
\end{equation}
All early fixed-interval integrals are finite, even if
$k/\rho-\ell_0$ initially has the opposite sign.
The matching is performed at each active receiver before
integration. Thus any cross contribution in excess of $O(\sqrt S)$
is matched by the negative mixed source. This does not assert that
the two have asymptotically equal integrals; the cross source
could itself be of order $\sqrt S$ or smaller.

\paragraph{Why these are the contributions on the initialized trajectory.}
Let $Q=\E_{P_2}(\xi^TD)^2$, $\alpha u+\delta v=
\E_{P_2}[(\xi^TD)r_t(D)]/Q$. SPM sources uses the overlap weight
$w_a=\E_{P_2}[(\xi^TD)^2{\bf1}_{a^TD>0}]/Q$.
On a probe-active label the original densities in question are
\[
 c_q=\tfrac Q{2q}\delta_+(v^Tb)^2w_a,\qquad
 n_q=\tfrac Q{2q}[-\alpha(v^Tb)(u^Tb)]_+w_a,\qquad
 p_q=\tfrac Q{2q}[\alpha(v^Tb)(u^Tb)]_+w_a,
\]
where $a,b$ are the current unit directions; all three are zero
on probe-inactive labels. Along resident extraction
$\alpha\to1$, $w_a\to P$ on a compact right chart and
$(u^Tb)/q\to\ell_0-y$. More precisely $\delta/q\to k/\rho-\ell_0$.
To see the constant, $v_1\to1$, $v_2/q\to\ell_0$ and the second
cluster has $\E D_1D_2=0$ and negligible fixed-time response.
The first-output fitted moment divided by $q$ gives $k/\rho$
by SPM divergence's Gaussian calculation. In using
$\xi^TD=-D_2+O(q|D_1|+q^2|D_2|)$ here, $D_1=O_{L^2(P_2)}(q)$;
the error divided by $q$ vanishes. Thus $c_q,n_q,p_q\to c_x,n_x,p_x$
on the active charts. Indicators may be passed under the integral:
different Q1 labels are strictly ordered in input at finite time
by their distinct remote incoming coefficients and the same-field
cooperative difference equation (the positive initial input gap
has a positive integrating-factor lower bound on a finite segment).
There can thus be no Q1 atom at a given input, since the label marginal
has no atoms; Q4 lies strictly below $\zeta$. Equivalently two
independent Q1 samples with distinct labels cannot tie.
Endpoint strips cost $O_S(\eta)$ for length $\eta$ since
$|W|=O_S(1)$, $|u^TW|=O_S(q)$ and $\delta/q=O_S(1)$ on the
fixed interval. The whole left mass is $o_S(q)$ by HBD--BRW.
Time integration and positive parts now pass by finite-horizon
domination. In particular \eqref{an04tbwv1:matched} bounds the
uncanceled part of SPM's very same divergent source.

\begin{theorem}[Trained surplus on an incoming interval]\label{an04tbwv1:surplus}
There exist moving pre-weak endpoints $b(q)\to\infty$ at which
the accumulated positive cross term diverges, the portion left
after receiver-by-receiver matching has square $O(b(q))$, and
the net longitudinal gain dominates the net transverse square.
In particular $A/\sqrt{b(q)}\to1$ and
\[
 \frac{A^2+q^2B^2}{2R_0B}\longrightarrow0,\qquad
 R_0=|r_{t_0}(\xi)|,
\]
at $t=t_h+b(q)$ for the actual $q$-flow. No conclusion at the
weak-response clock is included.
\end{theorem}
\begin{proof}
At fixed $S$, use SPM early to pass to
$A=H(S)-H(s_0)$ and
$B=\ell_0[H(S)-H(s_0)]-[V(S)-V(s_0)]$ in the limit,
$V=\int y(\zeta-x)_+\dd\mu$. In fact the COD core bound improves
to $V(S)\le-c e^{LS}/\sqrt{1+S}$ late: Q4 is included in the
upper face of Lemma~\ref{an04tbwv1:bulk}, hence its exposure is
at least $c\sqrt S$, its mass at least $1/4$, and its output at
most $-c'D(S)$. COD positive bounds the entire positive Q1 cost
by $o(D(S))$. Thus the limiting $B$ is at least
$c e^{LS}/\sqrt{1+S}$ for large $S$ uniformly over extractions.
SPM divergence's lower integrand $c(1+s)^{-5/6}$ gives an
unbounded lower integral for $c_q$ already in the iterated
resident limit. For integer $S=j\to\infty$ use compact extraction
on $[s_0,j]$ and the preceding bounds; choose $q_j\downarrow0$
so that for $q\le q_j$ the net transverse motion is
$\sqrt j+O(1)$, the unmatched positive integral is $O(\sqrt j)$,
the cumulative positive cross source exceeds a lower bound tending
to infinity, and $B\ge c e^{Lj}/\sqrt{1+j}$. Each upper bound
for $|B|$ is finite at fixed $j$, so also impose $q^2|B|<1/j$.
BRW delay permits $\sup_{0\le s\le j}|\mathsf m_2(t_h+s)|<1/j$.
Taking $b(q)=j$ on $(q_{j+1},q_j]$ proves all assertions. The
proof neither chooses nor controls the rate at which $b(q)$ grows.
\end{proof}

\section{An initialized seed for the attempt to reach weak learning}
\begin{theorem}[Fixed Q2 ancestry]\label{an04tbwv1:seed}
In physical balanced coordinates write $z=(U_2+W_2)/2$,
$H_2=\int_{\rm Q2} z^2\,dl$ with $dl=d\alpha/(2\pi)$.
At each fixed finite resident delay $s\ge0$,
\[
 \frac{H_2(t_h+s)}{q^2}\longrightarrow
                  \frac{c^{2\omega}}8(4e^s)^\lambda .
\]
This is a Cartesian amplitude, not a weak-chart mass count.
\end{theorem}
\begin{proof}
First observe at $T_A=\log(A/S_0)$ for fixed positive nominal $A$.
Fix a compact interior Q2 arc and a smaller fixed cap $a_0>0$.
At $T_{a_0}$ CRA outer gives the target angle tending to $\pi/2$
and $\widehat m=(1+o_q(1))S_0 e^{\lambda T_{a_0}}\sin^2\alpha$
uniformly on the arc since $\lambda/(1-\lambda)<1$.
P targetcomparison then bounds both trained angles within $O(a_0)$
of that angle and the relative radial weight by $e^{O(a_0)}$
times the target in both directions. In particular
\[
 \frac{z(T_{a_0})}{q}
   =c^\omega a_0^{\lambda/2}\sin\alpha\,
       (1+O(a_0)+o_{q;a_0}(1)).
\]
The big-$O$ is independent of the small cap. At fixed $a_0$,
both second components are bounded below by a positive multiple
of $q$, and $\sqrt m\le C_{a_0,A}q$ up to $T_A$ by the finite
elapsed time and HBD's generic radial bound.

For completeness, on this interval the full Gaussian fitted field
has $|f_1(D)|\le C_A(|D_1|+q|D_2|)$ and
$|f_2(D)|\le C_{a_0,A}(q|D_1|+q^2|D_2|)$, allowing $C_{a_0,A}$
also in the first bound. This is HBD--BRW physical on the fixed
nominal compact. For $R=\E[(D-f)D^T{\bf1}_{U^TD>0}]$ it implies
$|R_{12}|+|R_{21}|\le C_{a_0,A}q$ and $R_{22}\ge-C_{a_0,A}q^2$.
Both $U_2',W_2'\ge-C_{a_0,A}q^2$ while both components are
nonnegative. Integrating excludes a first contact of either with
half its positive lower entry bound, including joint contact.
Consequently $U_2/|U|$ is bounded away from zero there, uniformly
on the arc for fixed caps. Its $P_2$ wrong gate costs
$O_{a_0,A}(e^{-d_{a_0,A}/q^2})$ in the second diagonal data moment
(by a Gaussian tail and Cauchy--Schwarz). The $P_1$ second diagonal
data moment is $O(q^2)$ and the full fitted $22$ contribution is
$O_{a_0,A}(q^2)$. Hence $R_{22}=\lambda/2+O_{a_0,A}(q^2)$.
The exact equation
$z'=R_{22}z+(R_{12}W_1+R_{21}U_1)/2$ gives
$z'=\lambda z/2+O_{a_0,A}(q^2)$ on the arc.
Integrate to $T_A$: the cap-dependent accumulated error is $o(q)$
with caps fixed. The relative entry error $O(a_0)$ stays relative
after multiplication by $(A/a_0)^{\lambda/2}$. Take $q\to0$ and
then $a_0\to0$ to obtain $z(T_A)/q\to c^\omega A^{\lambda/2}\sin\alpha$.
These arguments hold uniformly for $A$ in a compact positive interval.
The gradient estimate \emph{does not} assert a bounded weak scaled
angle. HBD physical bounds $|z/q|$ on the entire Q2 circle at fixed
nominal times, so both endpoint strips have arbitrarily small
contribution after squaring. Integrate $\sin^2\alpha$ on Q2 against
$dl$ (integral $1/8$), and use $S_0e^{t_h+s}\to4e^s$.
\end{proof}

\begin{theorem}[Arrival at a fixed small weak response]\label{an04tbwv1:smallhit}
For the critical initialized Gaussian family there exist a fixed
$\eta_0>0$ and a fixed finite $b$ such that, for any fixed
$k_0>4/\lambda$ and all sufficiently small $q$, the first
post-strong-half hit of $\mathsf m_2=\eta_0$ occurs by
$t_h+b+k_0\log(1/q)$.
\end{theorem}
\begin{proof}
We first refine the tagged compact-Q2 entry at $t_b=t_h+b$.
There $U_1/q,W_1/q\to0$ uniformly on the arc. To check this
stronger statement, at the small nominal cap $a_0$ the upper
bound on these two absolute first components divided by $q$ is
$O(a_0^{\lambda/2+1})+o_{q;a_0}(1)$ by the same target comparison.
Thereafter put $D_1^*=\max(|U_1|,|W_1|)$. The off-diagonal
bounds and positive-component proof above give
\[
 D^+D_1^*\le a_*(1+S)D_1^*+C_{a_0,A}q^2,\quad
 a_*=1/2+q^2,\quad S=\int m\,dl.
\]
Here $|R_{11}|\le a_*(1+S)$ by the spectral bounds in P.
The homogeneous multiplier is bounded by
$C_A(A/a_0)^{a_*}$ since $\int_{T_{a_0}}^{T_A} S\,dt\le C_A$
from the global mass envelope, independent of $a_0$. The
inhomogeneous term is $o(q)$ at fixed caps. Thus the upper limit
after division by $q$ is $C_A a_0^{(1+\lambda)/2}$; send $a_0$
to zero. The moving clock again lies in a fixed compact of nominal
levels. Balance, the strictly positive second components, and the
seed formula on the arc now imply that $w=(U_2-W_2)/2=o(q)$ and
$z/q$ is bounded below uniformly there.

Write $r=D-f$. At any fixed $b$, HBD--BRW give
\[
 4\mathcal L(t_b)=\E_{P_1}|r|^2+\E_{P_2}|r|^2
     \longrightarrow (1-M(b))^2+\lambda .
\]
For this order-one statement, restrict first to compact right interiors:
their charted contributions converge in Gaussian $L^2$ to their
mass times $e_1(D_1)_+$ on $P_1$ and to zero on $P_2$.
Strips cost $O_b(\epsilon)$ in norm for label length $\epsilon$,
left ancestry has $o_b(1)$ mass and $D_1\to1$ on $P_1$.
Use the HBD mass clock. By loss monotonicity and the regression
inequality $\E_{P_2} r_2^2\ge(\lambda+q^2)(1-\mathsf m_2)^2$,
as long as $\mathsf m_2\le\eta_0<1$,
\begin{equation}\label{an04tbwv1:lossbudget}
 \E_{P_1}|r|^2+\E_{P_2}r_1^2
 \le (1-M(b))^2+\lambda(2\eta_0-\eta_0^2)+o_{q;b}(1).
\end{equation}
Choose $b$ large, then $\eta_0>0$ small so that the square root
of this upper bound is below $\lambda/400$ for small $q$.
Take also $\eta_0<1/8$. BRW ensures the weak response is below
this level up to $t_b$ for small $q$.

For all receiver gates, Cauchy--Schwarz on each cluster gives
$|R_{11}|,|R_{12}|,|R_{21}|\le\lambda/100$ throughout this
no-hit interval. Indeed each uses an $L^2$ residual component
bounded by \eqref{an04tbwv1:lossbudget}, except the $P_2$ term in
$R_{21}$, which is at most $q\|r_2\|_{P_2}/2=O(q)$.
All data component norms are at most two for small $q$. Whenever
the receiving input direction satisfies $U_2/|U|\ge1/\sqrt2$,
\[
 R_{22}\ge(\lambda+q^2)(1-\mathsf m_2)/2-O(q)
                    -O(e^{-d/q^2})\ge\lambda/3 .
\]
Indeed the full $P_2$ correlation is exact, its missing gate
moment is exponentially small by Cauchy--Schwarz, the Gaussian
tail on that cone and the bounded residual norm; the $P_1$ cost
is at most $q\|r_2\|_{P_1}/2$.

Start a positive-length compact Q2 arc at $t_b$. The reached entry
just proved is strict for the cone
$|U_1|,|W_1|,|w|\le z/2,\ z>0$.
Inside including its faces, $U_2,W_2\in[z/2,3z/2]$, hence
$U_2/|U|\ge1/\sqrt2$. The exact equations imply $z'\ge\lambda z/4$.
The upper derivatives of either absolute first component on its
face are at most $(\lambda/100)(z/2+3z/2)< z'/2$.
On $|w|=z/2$ the equation
$w'=-R_{22}w+(R_{12}W_1-R_{21}U_1)/2$ gives
$(|w|)'\le-\lambda z/6+\lambda z/200<z'/2$.
The common bounds cover simultaneous faces; the positive $z$ face
cannot be lost by the growth inequality. Thus the cone continues
for as long as there is no hit. On the arc
$\int z^2dl\ge c_b q^2 e^{\lambda(t-t_b)/2}$.
But $z^2\le m$, and $S\le C(1+t-t_h)$ by BRW, since
$S'=2\E[f\cdot(D-f)]\le\E|D|^2/2$. At duration
$k_0\log(1/q)$ with $k_0>4/\lambda$ these bounds contradict one
another. This proves the hit by continuity.
\end{proof}

The measured threshold here is small and fixed; no threshold
depending on $q$ was called learning. In particular the upper
bound in \eqref{an04tbwv1:lossbudget} at $\eta_0=1/4$
already includes $7\lambda/16$, greater than $(\lambda/400)^2$,
even before adding the entry loss. The sufficient small-coupling
test used on the cone is therefore unpaid at the required quarter
response, and still less so at the half hit. This is a failure of
that upper estimate, not a lower bound on the actual cross error.

\paragraph{Attempt to force a measured clock.}
Set $Q_2=\rho^2+q^2$, $m_2=\mathsf m_2$, and use the exact
CLOCK amplitude equation without division by $H_2$:
\begin{align*}
 H_2'&=Q_2(1-m_2)H_2+\mathcal D_2,\\
 \mathcal D_2&=\int_{\rm Q2}
 \{[2R_{22}-Q_2(1-m_2)]z^2+z(R_{12}W_1+R_{21}U_1)\}\,dl .
\end{align*}
All residual matrices include both sources and their actual gates.
With $G(t)=\exp[Q_2\int_{t_h}^t(1-m_2)]$, the normalized
retained amplitude is $H_2(t)/G(t)=H_2(t_h)+\int_{t_h}^t
\mathcal D_2/G$. Up to any time $T$ without a response hit
of $\eta\in(0,1)$ one therefore necessarily has
\begin{equation}\label{an04tbwv1:clockgap}
 H_2(t_h)+\int_{t_h}^T\mathcal D_2/G
 \le C(1+T-t_h)e^{-Q_2(1-\eta)(T-t_h)} .
\end{equation}
Indeed $H_2\le\int m$, and total radial mass is at most
$C(1+t-t_h)$ by the BRW physical mass identity.
At $T=t_h+k_0\log(1/q)$ with $k_0>2/[\lambda(1-\eta)]$,
the right side is $o(q^2)$. Theorem~\ref{an04tbwv1:seed} thus
pins down the charge: nonarrival by this time requires the
discounted defect to remove at least
$(c^{2\omega}4^\lambda/8-o(1))q^2$.
On each fixed delay interval $m_2=o(1)$ and the theorem gives
$\int\mathcal D_2/G=o(q^2)$ over that interval. This supplies no
such estimate on the logarithmic interval: the off-diagonal forcing
and receiver gate in $R_{22}$ are controlled above only with fixed
caps. For the required half response we cannot contradict
\eqref{an04tbwv1:clockgap} there. Even
a reached half-response hit would still require the learned-window
loss and alignment checks.

\section{Can feedback cancel the transverse source at the learned time?}
Here is a second attempt, valid without resident approximations and
even when the answer to the clock question is supplied. Follow the
very regression that produces the divergent term instead of discarding
its sign. In physical parameter space use the label $L^2(dl)$
inner product (both vector layers). Set
\[
 \ell=\nabla_\Theta\frac{\E_{P_2}[(\xi^TD)v^Tf_t(D)]}{Q},
 \qquad g_i=\nabla_\Theta[i^Tf_t(\xi)]\quad(i=u,v).
\]
In these formulas $u,v$ are the fixed witness axes, $Q=\E_{P_2}(\xi^TD)^2$.
The regression target is fixed, so $\delta'=-\langle\ell,\dot\Theta\rangle$.
Let $\mathcal N=\dot\Theta-Q\delta\ell/2$.
Equivalently $\mathcal N$ is the feature gradient applied to the
complete $P_1$ residual and the $P_2$ residual minus
$\delta(\xi^TD)v$, with both mixture weights. This is an explicit
field of the actual trajectory, including the weak longitudinal
response and the centered residual. Define
\[
 d=Q\|\ell\|^2/2,\quad p_i=Q\langle g_i,\ell\rangle/2,\quad
 F=-\langle\ell,\mathcal N\rangle,\qquad
 K_i(r,T)=\int_r^T p_i(t)\exp[-\textstyle\int_r^t d]\,dt .
\]
SPM's two-block product gives
$p_v=Q(j_{vv}+\omega_t)/2$ and $p_u=Qj_{uv}/2$; $\omega_t$ is
the direct-output term, including the actual mass, of SPM overlap.
With all signs retained the exact scalar and probe equations are
$\delta'=-d\delta+F$,
$i^T\dot f(\xi)=p_i\delta+\langle g_i,\mathcal N\rangle$.
Thus for any $t_b<T$ the remaining \emph{physical} probe movements are
\begin{equation}\label{an04tbwv1:filter}
 L_i(r,T):=i^T[f_T(\xi)-f_r(\xi)]
  =K_i(r,T)\delta(r)+
    \int_r^T\langle g_i(t)-K_i(t,T)\ell(t),\mathcal N(t)\rangle\,dt,
 \qquad r=t_b .
\end{equation}
Variation of constants and finite-interval Fubini give this identity;
it avoids differentiating a gate or dividing by a possibly small
instantaneous response rate. The explicit $\delta j_{vv}$ integral
has been eliminated, but its effect on the probe has not been set
to zero.

\paragraph{Test of the damping route on the actual entrance.}
It would be incorrect to treat $\delta$ as a freely decaying error.
Indeed for \emph{every} fixed $t_b=t_h+b$, $b\ge s_0$, and every
sequence of reached learned endpoints $T_q$,
\begin{equation}\label{an04tbwv1:false}
 \int_{t_b}^{T_q}\frac{F_+}{q}\,dt\longrightarrow\infty .
\end{equation}
On a fixed later resident interval $\delta(t_h+s)/q$ has limit
$k(s)/\rho-\ell_0\ge cD(s)-C_{s_0}$, while $|\delta(t_b)|/q\le C_b$.
Since $d\ge0$, the scalar equation bounds
$\delta(t_h+s)/q\le|\delta(t_b)|/q+\int_{t_b}^{t_h+s} F_+/q$.
First take small noise at fixed $s$, then send $s$ to infinity,
using BRW delay. Thus the shortcut of charging only a bounded
regeneration of the weak cross regression fails on the trained
family. This is not a sign claim for its \emph{filtered} net effect.

One can isolate the failure more sharply over the resident part.
Along extraction on fixed clocks, $d(t_h+s)\to
(\lambda/2)\int P^2 d\mu$, and $p_v(t_h+s)\to
(\lambda/2)\int_{\mathcal A}P\,d\mu$. For the first limit the input
block of $\ell$ has squared norm per radial mass tending to $P^2$
(the vector Gaussian moment tends to $-P e_2$); the output block
is $O(q)$ on bounded strong charts. On both blocks the squared
contribution per physical label is bounded by $C m$, so right strips
and negligible left mass can be removed as above. For $p_v$ use
the no-atom argument and the per-mass bound on $\omega_t$.
The protected outer labels have actual mass $\ge c/(1+s)$ with
positive input and $P\ge1/2$. Thus $d^{\lim}\ge c'/(1+s)$.
By \eqref{an04tbwv1:exposure}, $0\le p_v^{\lim}\le
C\log(2+s)/(1+s)$. For fixed $b$ it follows that
\[
 \sup_{S\ge b}\limsup_{q\downarrow0}|K_v(t_h+b,t_h+S)|
 \le C_b\int_b^\infty\log(2+s)(1+s)^{-1-c'}\,ds<\infty .
\]
This limit passage is on finite $[b,S]$ first. The homogeneous
response $K_v\delta(t_b)/q$ therefore cannot account for the
divergence of the resident signed cross integral. Indeed the latter
is eventually positive in the limit, with the SPM lower integral
of order at least $\int^S (1+s)^{-5/6}ds$; the direct-output
correction $\omega_t\delta/q$ vanishes on each fixed interval.
The continuing forcing of the regression is essential. This is
consistent with the cancellation estimate \eqref{an04tbwv1:matched}.

At times where the exact projected weak response
$h=\E_{P_2}[(\xi^TD)u^Tf_t(D)]/Q$ has $|h|\ge\eta'>0$,
$\|\ell\|^2\ge(\eta')^2/\int (u^TW)^2 dl$ by the output
block and Cauchy--Schwarz. SPM shows $h=\mathsf m_2+O_{s_0}(q)$
uniformly from the loss bound. Thus there is damping at a learned
time. Even on an entire interval with this guard it controls the
kernel $K_i$ only by a weighted average of $p_i/d$; for instance
$|K_i|\le C\sup(\int m\,dl)^2/(\eta')^2$ on that interval
by the Gaussian moment bounds. The guard is absent on most of
the incoming interval; the resulting bound is unsigned and does not
control the correlations against $\mathcal N$ at the required scales.

\paragraph{Exact learned-endpoint deficit.}
Choose a fixed late $b$ with the incoming surplus above, put
$A_b=A(t_b)$, $B_b=B(t_b)$, $U_b=R_0-q^2B_b$.
Let $T$ be in the actual measured learned window of DRR (if reached).
The exact test on the remaining passage is
\begin{equation}\label{an04tbwv1:missing}
 \frac{A_b L_v(t_b,T)}q+\frac{L_v(t_b,T)^2}{2q^2}
 -\frac{U_b L_u(t_b,T)}{q^2}+\frac{L_u(t_b,T)^2}{2q^2}
       < e_*-e_q(t_b).
\end{equation}
This is obtained by subtracting the DRR squared-residual identity
at $t_b$ from the one at $T$; the four terms include the interaction
with the transverse displacement already acquired. All quantities
on the left are specified regression-feedback integrals
\eqref{an04tbwv1:filter}. The incoming right side is at least of
order $e^{Lb}/\sqrt{1+b}$ for late fixed $b$ in the small-noise
limit. That lower bound depends on choosing $q$ after $b$.
For a learned-depth theorem one needs the strict inequality at
\emph{some} reached learned $T$; a fixed asymptotic margin yields
a quantitative gain. We do not have a signed joint bound on the
correlations $\langle g_i-K_i\ell,\mathcal N\rangle$ there,
or a continuation of the resident activity and damping estimates
to this time. In particular \eqref{an04tbwv1:false} rules out the
bounded positive-forcing shortcut but does not refute
\eqref{an04tbwv1:missing}. Replacing $b$ by the learned delay in
the fixed-clock bounds would assume the missing passage.

\paragraph{Dependency ledger.}
HBD physical, outer, incoming and clock and BRW budget, delay and
field apply to this critical family on fixed nominal or resident
intervals. BRW's protected-set conclusions use the HBD entrance
and logistic mass, including its radial mass factor. COD mean,
exposure and positive concern that same limiting law and supply
fixed-time probe transfer. SPM early, sources, overlap and divergence
use the fixed critical witness axes. The Q2 seed uses CRA outer's
positive-noise target comparison on compact interior Q2 labels
at exponent $1/\omega$, P targetcomparison only before the small cap,
and the COD compact-Q2 finite-time physical bounds thereafter.
CLOCK's Cartesian identity is proved for fixed ancestry irrespective
of gate or chart; its conditional deep passage estimates are not
imported. No generated trajectory or finite-noise empirical value
enters the argument.
\end{document}
